% Hand-written sentences that depend on results; every number is a macro from v16_numbers.tex (Research/code/v16_tables.py).
\newcommand{\nAbstractModel}{a log-linear working model fitted to the same capped records was more accurate than the design-based score for every tree up to \nCurveLastAllBetter{} records, and less accurate for a third of the trees at 16,000 and beyond, where its bias no longer shrinks.}
\newcommand{\nPointerSentence}{at $b\in\{3,4\}$ fixed forms lose exactly $4$ and $8$, two rounds lose $1.80$ and $3.32$--$3.33$ against the upper bounds $2$ and $4$ of part (a), and $b-1$ rounds lose $1.80$ at both caps (Appendix~\ref{app:pointerexact}).}
\newcommand{\nDiscussionSentence}{optimized forms needed between $1.0$ and \nMaxAll{} times the records of branching; one adaptive round recovered at least \nTwoMinPct\% of the log-gap wherever its optimum was computed; label-dependent form probabilities did not change the picture; and learned designs remained exactly unbiased at an all-in price within \nAllInOverFourd--\nAllInOverFived\% of the oracle for audits of 8,000 records at caps four and five.}
\newcommand{\nOrderSentence}{a 10\% chance that the last answer flips biases branching and form scores by a median \nPosBranch{} and \nPosFit{} branching standard errors at $n=1000$ (Section~\ref{sec:experiments}).}
\newcommand{\nModelParagraph}{The constants above are for design-unbiased estimation; a model can buy accuracy, at a price. A log-linear working model for the label and the rule's items, fitted by EM to the records of the fitted forms, needs no reference \citep{schafer1997}. Figure~\ref{fig:rounds}(b) tracks its RMSE against the design-based form score on the same records for \nCurveTrees{} trees. With two-way interactions the model is the more accurate for every tree up to $n=\nCurveLastAllBetter$ (median RMSE ratio \nCurveTwoMedLo--\nCurveTwoMedHi{} at every size, on a par with branching at \nCurveBranchMed). Its bias does not shrink: at $n=16{,}000$ and $32{,}000$ it is the less accurate for \nCurveTwoWorsec{} and \nCurveTwoWorsed{} of the \nCurveTrees{} trees, by up to a factor \nCurveTwoMaxc{} and \nCurveTwoMaxd, and its jackknife intervals cover \nEmBigCovLo--\nEmBigCovHi{} at $n=16{,}000$ where the design-based intervals keep \nEmBigFixCovLo--\nEmBigFixCovHi{} (Table~\ref{tab:em}). With three-way interactions the model stays within \nCurveThreeLo--\nCurveThreeHi{} of the form score at every size. Imputation from complete records of another population is more accurate still when the conditional law of the unasked items transports (RMSE \nImpRmseLo--\nImpRmseHi{} times that of branching), and covers only \nImpCovLo--\nImpCovHi{} under \nImpBadN{} of eight pre-specified shift settings (Table~\ref{tab:shift}). At survey sizes of a few thousand records, a parsimonious model on optimized forms is therefore the more accurate tool, and adaptivity adds little to it. The constants of this paper price pre-commitment for claims that must not rest on a model, and for large audits, where the model fails for a minority of rules that the capped records do not identify.}
\newcommand{\nEmAppendixText}{At $n=1000$ the working model on the fitted-form records has the smallest RMSE of the four estimators in most settings, \nEmBrLo--\nEmBrHi{} times that of branching: the model shares strength across paths that the design-based scores treat separately. Its bias is between \nEmBiasLo{} and \nEmBiasHi{} branching standard deviations and its coverage \nEmCovLo--\nEmCovHi{} (below 0.925 in \nEmBelow{} of \nEmSettings{} settings). The bias is a property of the model, not of the sample size: the accuracy of a depth-$b$ tree depends on interactions of up to $b$ items and the label, which a two-way model does not represent. At $n=\nEmBigN$ the same bias is up to \nEmBigBiasHi{} of the estimator\textquotesingle s own standard deviations, coverage falls to \nEmBigCovLo--\nEmBigCovHi{} (below 0.90 in \nEmBigBelow{} of \nEmBigSettings{} settings), and the RMSE is \nEmBigRmseLo--\nEmBigRmseHi{} times that of branching and above that of the design-based form score in \nEmBigWorse{} settings; the design-based scores keep coverage \nEmBigFixCovLo--\nEmBigFixCovHi. A richer model would reduce the bias and raise the variance; with all interactions it is the design-based estimator. On balanced random forms the model-based estimate is unbiased to a similar degree but \nEmUnifLo--\nEmUnifHi{} times less accurate at $n=1000$ than the design-based score on optimized forms, so the choice of forms matters for model-based analysis as well.}
